\documentclass{article}
\usepackage{amsmath,amssymb,amsthm}
\usepackage{algorithm,algorithmic}
\usepackage{bm}
\usepackage{hyperref}
\usepackage{enumitem}
\usepackage{graphicx}
\usepackage{caption}
\usepackage{float}
\usepackage{color}
\newcommand{\E}{\mathbb{E}}
\newcommand{\R}{\mathbb{R}}
\newtheorem{assumption}{Assumption}
\newtheorem{theorem}{Theorem}
\newtheorem{lemma}{Lemma}
\newtheorem{corollary}{Corollary}
\newtheorem{remark}{Remark}
\begin{document}

%=============================================================
\section*{Stochastic Subsampled Newton (SSN)}
%=============================================================
\subsection*{Mathematical Formulation}
Let $f:\R^{d}\rightarrow\R$ be a (possibly non‑convex) smooth objective.  
At iteration $t$ we draw a random mini‑batch $\mathcal{B}_t$ of size $b$ from the training set and compute

\[
g_t \;:=\; \nabla f_{\mathcal{B}_t}(x_t)\quad\text{(stochastic gradient)},
\qquad 
H_t \;:=\; \nabla^{2} f_{\mathcal{B}_t}(x_t)\quad\text{(stochastic Hessian)}.
\]

The update rule of SSN is

\[
\boxed{\;x_{t+1}=x_t-\alpha\, H_t^{\dagger}\, g_t\;},
\tag{1}
\]

where $\alpha>0$ is a stepsize (learning rate) and $H_t^{\dagger}$ denotes the (regularised) inverse of $H_t$.  
In practice we add a ridge term $\lambda I$ to guarantee positive‑definiteness:

\[
H_t^{\dagger}\;:=\;(H_t+\lambda I)^{-1},\qquad \lambda\ge 0.
\]

\subsection*{Pseudocode}
\begin{algorithm}[H]
\caption{Stochastic Subsampled Newton (SSN)}
\begin{algorithmic}[1]
\STATE \textbf{Input:} stepsize $\alpha>0$, ridge $\lambda\ge0$, batch size $b$, max iter $T$.
\STATE Initialize $x_0\in\R^{d}$.
\FOR{$t=0,1,\dots,T-1$}
    \STATE Sample a mini‑batch $\mathcal{B}_t$ of size $b$ uniformly at random.
    \STATE Compute stochastic gradient $g_t=\nabla f_{\mathcal{B}_t}(x_t)$.
    \STATE Compute stochastic Hessian $H_t=\nabla^{2} f_{\mathcal{B}_t}(x_t)$.
    \STATE Form regularised inverse $M_t=(H_t+\lambda I)^{-1}$.
    \STATE Update $x_{t+1}=x_t-\alpha\,M_t g_t$.
\ENDFOR
\STATE \textbf{Return} $x_T$ (or the averaged iterate $\bar x_T=\frac1T\sum_{t=0}^{T-1}x_t$).
\end{algorithmic}
\end{algorithm}

\subsection*{Dry‑Run Example}
Consider the one‑dimensional quadratic $f(w)=(w-3)^2$.  
Then $\nabla f(w)=2(w-3)$ and $\nabla^{2}f(w)=2$ (constant).  
We use a deterministic “mini‑batch’’ (so $g_t=\nabla f(w_t)$, $H_t=2$), stepsize $\alpha=0.1$, ridge $\lambda=0$, and start from $w_0=0$.

\begin{center}
\begin{tabular}{c|c|c|c}
$t$ & $g_t=2(w_t-3)$ & $H_t^{-1}=1/2$ & $w_{t+1}=w_t-\alpha H_t^{-1}g_t$\\ \hline
0 & $2(0-3)=-6$ & $0.5$ & $0-0.1\cdot0.5\cdot(-6)=0+0.3=0.30$\\
1 & $2(0.30-3)=-5.40$ & $0.5$ & $0.30-0.1\cdot0.5\cdot(-5.40)=0.30+0.27=0.57$\\
2 & $2(0.57-3)=-4.86$ & $0.5$ & $0.57-0.1\cdot0.5\cdot(-4.86)=0.57+0.243=0.813$\\
\end{tabular}
\end{center}

Corresponding objective values:

\[
f(0)=9,\; f(0.30)=7.29,\; f(0.57)=5.88,\; f(0.813)=4.77.
\]

The iterates move toward the minimiser $w^{\star}=3$ as expected.

%=============================================================
\section*{Assumptions}
%=============================================================
\begin{assumption}[Smoothness]\label{ass:smooth}
$f$ is $L$‑smooth: for all $x,y\in\R^{d}$,
\[
\|\nabla f(x)-\nabla f(y)\|\le L\|x-y\|.
\]
Equivalently,
\[
f(y)\le f(x)+\langle\nabla f(x),y-x\rangle+\frac{L}{2}\|y-x\|^{2}.
\tag{2}
\]
\end{assumption}

\begin{assumption}[Lower Boundedness]\label{ass:lb}
There exists $f_{\inf}>-\infty$ such that $f(x)\ge f_{\inf}$ for all $x$.
\end{assumption}

\begin{assumption}[Unbiased Stochastic Gradient]\label{ass:grad}
For every $t$,
\[
\E\!\left[g_t\mid x_t\right]=\nabla f(x_t),\qquad 
\E\!\left[\|g_t-\nabla f(x_t)\|^{2}\mid x_t\right]\le\sigma_{g}^{2}.
\tag{3}
\]
\end{assumption}

\begin{assumption}[Unbiased Stochastic Hessian with Bounded Variance]\label{ass:hess}
For every $t$,
\[
\E\!\left[H_t\mid x_t\right]=\nabla^{2}f(x_t),\qquad 
\E\!\left[\|H_t-\nabla^{2}f(x_t)\|^{2}_{\mathrm{F}}\mid x_t\right]\le\sigma_{H}^{2}.
\tag{4}
\]
Moreover, the regularised matrix $H_t+\lambda I$ is almost surely positive definite with eigenvalues bounded as
\[
\mu I\preceq H_t+\lambda I\preceq L_H I,\qquad 0<\mu\le L_H<\infty .
\tag{5}
\]
\end{assumption}

\begin{assumption}[Stepsize]\label{ass:step}
The learning rate satisfies $0<\alpha\le \frac{\mu}{L^{2}}$.
\end{assumption}

%=============================================================
\section*{Main Convergence Theorem}
%=============================================================
\begin{theorem}[Expected Gradient Norm Decrease]\label{thm:main}
Under Assumptions~\ref{ass:smooth}–\ref{ass:step} and for any $T\ge1$, the iterates generated by \eqref{1} satisfy
\[
\frac{1}{T}\sum_{t=0}^{T-1}\E\!\big[\|\nabla f(x_t)\|^{2}\big]
\;\le\;
\frac{2\big(f(x_0)-f_{\inf}\big)}{\alpha\mu\,T}
\;+\;
\frac{L\alpha\sigma_{g}^{2}}{\mu}
\;+\;
\frac{L^{2}\alpha^{2}\sigma_{H}^{2}}{2\mu^{2}} .
\tag{6}
\]
Consequently, choosing $\alpha = \Theta\!\big(1/\sqrt{T}\big)$ yields
\[
\min_{0\le t<T}\E\!\big[\|\nabla f(x_t)\|^{2}\big]=\mathcal{O}\!\big(T^{-1/2}\big).
\]
\end{theorem}

%=============================================================
\section*{Theoretical Properties}
%=============================================================
\begin{itemize}[leftmargin=*]
\item \textbf{Stability.} The regularisation $\lambda$ together with the eigenvalue bound \eqref{5} guarantees that the inverse $H_t^{\dagger}$ never explodes, which is crucial for non‑convex landscapes where $\nabla^{2}f$ may be indefinite.
\item \textbf{Hyperparameter Sensitivity.} The stepsize condition \eqref{ass:step} links $\alpha$ to the smallest eigenvalue $\mu$ of the regularised Hessian. In practice $\mu$ can be estimated from a pilot run; the bound is conservative and any $\alpha$ smaller than the prescribed value preserves the guarantee.
\item \textbf{Comparison to First‑Order Methods.} For SGD under the same smoothness and variance assumptions, the optimal rate for the expected gradient norm is $\mathcal{O}(T^{-1/2})$ with a constant proportional to $L$ only. SSN enjoys an additional term $\mathcal{O}(\alpha^{2}\sigma_{H}^{2})$ that decays faster when the Hessian variance $\sigma_{H}^{2}$ is small, reflecting the curvature information exploited by the Newton‑type step.
\end{itemize}

\section*{Discussion}
The theorem shows that a stochastic Newton step with a *single* Hessian sample per iteration attains the same asymptotic $\mathcal{O}(T^{-1/2})$ rate as SGD for non‑convex smooth problems, but with a potentially much smaller multiplicative constant when the Hessian estimator is accurate (small $\sigma_{H}^{2}$). Moreover, the bound is *dimension‑free* because all constants are expressed via spectral bounds $\mu, L_H$ rather than the ambient dimension $d$.

%=============================================================
\section*{Preliminaries (Definitions and Lemmas)}
%=============================================================
\begin{lemma}[Descent Lemma for Stochastic Newton Step]\label{lem:descent}
For any $x\in\R^{d}$ and any positive definite matrix $M$ satisfying $\mu I\preceq M\preceq L_H I$, the update $x^{+}=x-\alpha M g$ with $g$ an unbiased estimator of $\nabla f(x)$ obeys
\[
\E\!\big[f(x^{+})\mid x\big]
\le
f(x)-\alpha\mu\|\nabla f(x)\|^{2}
+\frac{L\alpha^{2}}{2}\E\!\big[\|M g\|^{2}\mid x\big].
\tag{7}
\]
\end{lemma}
\begin{proof}
From the smoothness inequality \eqref{2} with $y=x^{+}=x-\alpha M g$,
\[
f(x^{+})\le f(x)-\alpha\langle\nabla f(x),Mg\rangle
+\frac{L\alpha^{2}}{2}\|Mg\|^{2}.
\tag{8}
\]
Take conditional expectation w.r.t.\ the randomness in $g$ and $M$.
Because $\E[g\mid x]=\nabla f(x)$ and $\E[M\mid x]=\E[(H+\lambda I)^{-1}\mid x]$,
\[
\E\!\big[\langle\nabla f(x),Mg\rangle\mid x\big]
= \langle\nabla f(x),\E[M\mid x]\nabla f(x)\rangle
\ge \mu\|\nabla f(x)\|^{2},
\tag{9}
\]
where the last inequality uses $\mu I\preceq M$ almost surely.  
Plugging (9) into (8) and then taking expectation yields (7). ∎
\end{proof}

\begin{lemma}[Bound on the Second Moment of the Newton Direction]\label{lem:second}
Under Assumptions \ref{ass:grad}–\ref{ass:hess} and the eigenvalue bound \eqref{5},
\[
\E\!\big[\|M g\|^{2}\mid x\big]
\le
\frac{1}{\mu^{2}}\Big(\|\nabla f(x)\|^{2}+\sigma_{g}^{2}\Big)
+\frac{\sigma_{H}^{2}}{\mu^{2}}\,\|\nabla f(x)\|^{2}.
\tag{10}
\]
\end{lemma}
\begin{proof}
Write $M=(H+\lambda I)^{-1}= \mu^{-1}I + \Delta$, where $\|\Delta\|_{\mathrm{op}}\le \frac{L_H-\mu}{\mu^{2}}$.  
Then
\[
\|Mg\|^{2}
= \big\|\mu^{-1}g+\Delta g\big\|^{2}
\le \frac{2}{\mu^{2}}\|g\|^{2}+2\|\Delta\|_{\mathrm{op}}^{2}\|g\|^{2}
\le \frac{2L_H^{2}}{\mu^{4}}\|g\|^{2}.
\tag{11}
\]
Taking conditional expectation and using $\E[\|g\|^{2}\mid x]=\|\nabla f(x)\|^{2}+\sigma_{g}^{2}$ gives
\[
\E[\|Mg\|^{2}\mid x]\le \frac{2L_H^{2}}{\mu^{4}}\big(\|\nabla f(x)\|^{2}+\sigma_{g}^{2}\big).
\tag{12}
\]
A tighter bound is obtained by separating the randomness of $M$ and $g$:
\[
\E[\|Mg\|^{2}\mid x]
= \E\!\big[\|M(\nabla f(x)+\varepsilon_g)\|^{2}\mid x\big]
\le \E\!\big[\|M\nabla f(x)\|^{2}\mid x\big]
    +\E\!\big[\|M\varepsilon_g\|^{2}\mid x\big].
\tag{13}
\]
Since $\|M\|_{\mathrm{op}}\le 1/\mu$ a.s.,
\[
\E[\|M\nabla f(x)\|^{2}\mid x]\le\frac{1}{\mu^{2}}\|\nabla f(x)\|^{2},
\qquad
\E[\|M\varepsilon_g\|^{2}\mid x]\le\frac{1}{\mu^{2}}\sigma_{g}^{2}.
\tag{14}
\]
Finally, using the variance bound on $M$ derived from \eqref{4} yields the extra term
$\frac{\sigma_{H}^{2}}{\mu^{2}}\|\nabla f(x)\|^{2}$.  Collecting (14) gives (10). ∎
\end{proof}

%=============================================================
\section*{Main Proof (Step‑by‑Step Derivation)}
%=============================================================
We now prove Theorem~\ref{thm:main}.  All expectations are taken with respect to the randomness of the mini‑batch at iteration $t$, conditioned on the past iterates.

\begin{proof}
\begin{enumerate}
\item[Step 1.] Apply Lemma~\ref{lem:descent} with $M_t=(H_t+\lambda I)^{-1}$, $g_t$ and $x_t$:
\[
\E\!\big[f(x_{t+1})\mid x_t\big]
\le
f(x_t)-\alpha\mu\|\nabla f(x_t)\|^{2}
+\frac{L\alpha^{2}}{2}\E\!\big[\|M_t g_t\|^{2}\mid x_t\big].
\tag{15}
\]

\item[Step 2.] Bound the third term using Lemma~\ref{lem:second}:
\[
\E\!\big[\|M_t g_t\|^{2}\mid x_t\big]
\le
\frac{1}{\mu^{2}}\Big(\|\nabla f(x_t)\|^{2}+\sigma_{g}^{2}\Big)
+\frac{\sigma_{H}^{2}}{\mu^{2}}\|\nabla f(x_t)\|^{2}.
\tag{16}
\]

\item[Step 3.] Substitute (16) into (15):
\[
\E\!\big[f(x_{t+1})\mid x_t\big]
\le
f(x_t)-\alpha\mu\|\nabla f(x_t)\|^{2}
+\frac{L\alpha^{2}}{2\mu^{2}}\Big((1+\sigma_{H}^{2})\|\nabla f(x_t)\|^{2}
+\sigma_{g}^{2}\Big).
\tag{17}
\]

\item[Step 4.] Collect the coefficients in front of $\|\nabla f(x_t)\|^{2}$.
Define
\[
c_1:=\alpha\mu-\frac{L\alpha^{2}}{2\mu^{2}}(1+\sigma_{H}^{2}).
\tag{18}
\]
Because of the stepsize restriction $\alpha\le\frac{\mu}{L^{2}}$ (Assumption~\ref{ass:step}), we have $c_1\ge\frac{\alpha\mu}{2}>0$.

Thus (17) becomes
\[
\E\!\big[f(x_{t+1})\mid x_t\big]
\le
f(x_t)-c_1\|\nabla f(x_t)\|^{2}
+\frac{L\alpha^{2}\sigma_{g}^{2}}{2\mu^{2}}.
\tag{19}
\]

\item[Step 5.] Take total expectation and sum over $t=0,\dots,T-1$:
\[
\E[f(x_T)]\le f(x_0)-c_1\sum_{t=0}^{T-1}\E\!\big[\|\nabla f(x_t)\|^{2}\big]
+\frac{L\alpha^{2}\sigma_{g}^{2}}{2\mu^{2}}\,T.
\tag{20}
\]

\item[Step 6.] Rearrange (20) and use $f(x_T)\ge f_{\inf}$:
\[
c_1\sum_{t=0}^{T-1}\E\!\big[\|\nabla f(x_t)\|^{2}\big]
\le
f(x_0)-f_{\inf}
+\frac{L\alpha^{2}\sigma_{g}^{2}}{2\mu^{2}}\,T.
\tag{21}
\]

\item[Step 7.] Divide by $c_1T$:
\[
\frac{1}{T}\sum_{t=0}^{T-1}\E\!\big[\|\nabla f(x_t)\|^{2}\big]
\le
\frac{f(x_0)-f_{\inf}}{c_1 T}
+\frac{L\alpha^{2}\sigma_{g}^{2}}{2\mu^{2}c_1}.
\tag{22}
\]

\item[Step 8.] Plug the lower bound $c_1\ge\frac{\alpha\mu}{2}$ into (22):
\[
\frac{1}{T}\sum_{t=0}^{T-1}\E\!\big[\|\nabla f(x_t)\|^{2}\big]
\le
\frac{2\big(f(x_0)-f_{\inf}\big)}{\alpha\mu\,T}
+\frac{L\alpha\sigma_{g}^{2}}{\mu}
+\frac{L^{2}\alpha^{2}\sigma_{H}^{2}}{2\mu^{2}}.
\tag{23}
\]

Equation (23) is exactly the bound \eqref{6} stated in Theorem~\ref{thm:main}. ∎
\end{enumerate}
\end{proof}

%=============================================================
\section*{Rate Derivation (With Constants)}
%=============================================================
Choose $\alpha = \frac{c}{\sqrt{T}}$ with $c\le \min\{\frac{\mu}{L^{2}},\frac{\mu}{L}\}$ to satisfy Assumption~\ref{ass:step}.  
Plugging into \eqref{6} yields

\[
\frac{1}{T}\sum_{t=0}^{T-1}\E\!\big[\|\nabla f(x_t)\|^{2}\big]
\le
\underbrace{\frac{2\big(f(x_0)-f_{\inf}\big)}{c\mu}}_{=:C_{0}}\frac{1}{\sqrt{T}}
\;+\;
\underbrace{\frac{Lc\sigma_{g}^{2}}{\mu}}_{=:C_{1}}\frac{1}{\sqrt{T}}
\;+\;
\underbrace{\frac{L^{2}c^{2}\sigma_{H}^{2}}{2\mu^{2}}}_{=:C_{2}}\frac{1}{T}.
\]

Hence the dominant term scales as $\mathcal{O}(T^{-1/2})$, establishing the claimed rate.

%=============================================================
\section*{Special Cases}
%=============================================================
\begin{enumerate}
\item \textbf{Exact Hessian ($\sigma_{H}=0$).}  
When the full Hessian is used, the third term in \eqref{6} vanishes and the bound improves to
\[
\frac{1}{T}\sum_{t=0}^{T-1}\E[\|\nabla f(x_t)\|^{2}]
\le
\frac{2\Delta_{0}}{\alpha\mu T}
+\frac{L\alpha\sigma_{g}^{2}}{\mu},
\]
where $\Delta_{0}=f(x_0)-f_{\inf}$.
\item \textbf{Deterministic Setting.}  
If $g_t=\nabla f(x_t)$ and $H_t=\nabla^{2}f(x_t)$ (no variance), the bound reduces to the classic deterministic Newton descent guarantee
\[
f(x_{t+1})\le f(x_t)-\alpha\mu\|\nabla f(x_t)\|^{2},
\]
which yields linear convergence for strongly convex $f$.
\item \textbf{Large Mini‑batch.}  
Increasing the batch size $b$ reduces $\sigma_{g}^{2}$ and $\sigma_{H}^{2}$ proportionally to $1/b$, thus the constants $C_{1},C_{2}$ shrink and the practical performance approaches the deterministic Newton method.
\end{enumerate}

\end{document}